\documentclass[a4paper]{arrowhead}

\usepackage[yyyymmdd]{datetime}
\usepackage{etoolbox}
\usepackage[utf8]{inputenc}
\usepackage[T1]{fontenc}
\usepackage{multirow}
\usepackage{longtable}
\usepackage{tikz}
\usetikzlibrary{positioning,shapes.multipart,arrows.meta}

\renewcommand{\dateseparator}{-}
\setlength{\emergencystretch}{3em}
\lstset{language=HTTP,breakatwhitespace=false,columns=fullflexible}

%% Special references
\newcommand{\mref}[2][]{{\textcolor{ArrowheadPurple}{\hyperref[sec:model:#2]{\ifstrempty{#1}{#2}{#1}}}}}
\newcommand{\oref}[1]{{\textcolor{ArrowheadBlue}{\hyperref[sec:operations:#1]{#1}}}}
\newcommand{\pdef}[1]{{\textcolor{ArrowheadGrey}{#1\label{sec:model:primitives:#1}\label{sec:model:primitives:#1s}\label{sec:model:primitives:#1es}}}}
\newcommand{\pref}[1]{{\textcolor{ArrowheadGrey}{\hyperref[sec:model:primitives:#1]{#1}}}}

% as template_sd.tex; the input and output arguments are passed already formatted (\mref or \pref)
\newrobustcmd\osubsection[3]{
  \addtocounter{subsection}{1}
  \addcontentsline{toc}{subsection}{\protect\numberline{\thesubsection}operation \textcolor{ArrowheadBlue}{#1}}
  \subsection*{
    \thesubsection\quad
    operation
    \textcolor{ArrowheadBlue}{#1}
    (\notblank{#2}{#2}{})
    \notblank{#3}{: #3}{}
  }
  \label{sec:operations:#1}
}
\newrobustcmd\msubsection[2]{
  \addtocounter{subsection}{1}
  \addcontentsline{toc}{subsection}{\protect\numberline{\thesubsection}#1 \textcolor{ArrowheadPurple}{#2}}
  \subsection*{\thesubsection\quad#1 \textcolor{ArrowheadPurple}{#2}}
  \label{sec:model:#2} \label{sec:model:#2s} \label{sec:model:#2es}
}
%%

\begin{document}
\sloppy

%% Arrowhead Document Properties
\ArrowheadTitle{gateTunnelConnection}
\ArrowheadServiceID{gateTunnelConnection}
\ArrowheadType{Service Description}
\ArrowheadTypeShort{SD}
\ArrowheadVersion{5.2.2}
\ArrowheadDate{2026-10-30}
\ArrowheadAuthor{Jerker Delsing}
\ArrowheadStatus{RELEASE CANDIDATE}
\ArrowheadContact{jerker.delsing@ltu.se}
\ArrowheadFooter{\href{www.arrowhead.eu}{www.arrowhead.eu}}
\ArrowheadSetup
%%

%% Front Page
\begin{center}
  \vspace*{1cm}
  \huge{\arrowtitle}

  \vspace*{0.2cm}
  \LARGE{\arrowtype}
  \vspace*{1cm}

  \Large{Service ID: \textit{"\arrowid"}}

  \vspace*{\fill}

  % Front Page Image

  \vspace*{1cm}
  \vspace*{\fill}

  % Front Page Abstract
  \begin{abstract}
    This document describes the \texttt{gate\allowbreak{}Tunnel\allowbreak{}Connection} service, version 5.2.2, produced by the GateTunnel support system of Arrowhead 5.2.2.
  \end{abstract}

  \vspace*{1cm}
\end{center}

\newpage
%%

%% Table of Contents
\tableofcontents
\newpage
%%

\section{Overview}
\label{sec:overview}

This document describes the \texttt{gate\allowbreak{}Tunnel\allowbreak{}Connection} service, which provides relay sessions between the local cloud and a neighbour cloud: provider and consumer sessions opened by the Gatekeeper, and the legacy control sessions.
It is produced by the GateTunnel system \cite{gt-sysd} in service version 5.2.2.

The rest of this document is organized as follows.
In Section \ref{sec:operations}, we describe the abstract message operations provided by the service.
In Section \ref{sec:model}, we end the document by presenting the data types used by the mentioned operations.
The concrete interface is specified in the HTTP/TLS/JSON IDD \cite{gt-gateTunnelConnection-http-idd}.

\subsection{Significant Prior Art}

Successor of the Arrowhead 4.6 Gateway (\texttt{/\allowbreak{}gateway} paths are still accepted).

\subsection{How This Service Is Meant to Be Used}

The local Gatekeeper opens a provider session (\texttt{open-\allowbreak{}provider-\allowbreak{}session}) when a neighbour's consumer negotiated one of our providers, and a consumer session (\texttt{open-\allowbreak{}consumer-\allowbreak{}session}) for a local consumer; the consumer session listens on a local port, and every request is forwarded over the relay. \texttt{connect-\allowbreak{}provider}, \texttt{connect-\allowbreak{}consumer} and \texttt{close-\allowbreak{}sessions} keep the v4 control-session API.

\subsection{Important Delimitations and Dependencies}
\label{sec:delimitations}

Only the Gatekeeper opens relay sessions. Sessions are restricted to the negotiated peer, consumer and methods/paths; default lifetimes 12 h maximum and 1 h idle. Consumer sessions are volatile and renegotiated after restart. The legacy control sessions have no data plane.

\newpage

\section{Service Interface}
\label{sec:operations}

This section describes the interfaces to the \texttt{gate\allowbreak{}Tunnel\allowbreak{}Connection} service.
In particular, each subsection names an abstract operation, an input type and an output type, in that order.
The input type is named inside parentheses, while the output type is preceded by a colon.
Input and output types are only denoted when accepted or returned, respectively, by the interface in question.

All abstract data types named in this section are defined in Section \ref{sec:model}.

\begin{figure}[ht!]
  \centering
  \begin{tikzpicture}[font=\small]
    \node[draw=ArrowheadBlue, very thick, rectangle split, rectangle split parts=2, align=center,
          rectangle split part fill={ArrowheadBlue!12,white}, text width=9cm, inner sep=6pt] (b)
      {\guillemotleft{}Service\guillemotright{}\\\textbf{gateTunnelConnection}
        \nodepart{second} \texttt{connect-\allowbreak{}provider}\\ \texttt{connect-\allowbreak{}consumer}\\ \texttt{close-\allowbreak{}sessions}\\ \texttt{open-\allowbreak{}provider-\allowbreak{}session}\\ \texttt{open-\allowbreak{}consumer-\allowbreak{}session}};
  \end{tikzpicture}
  \caption{SysML-style block of the \texttt{gate\allowbreak{}Tunnel\allowbreak{}Connection} service and its operations.}
  \label{fig:SD}
\end{figure}

The following interface operations are available.

\osubsection{connect-provider}{\mref{ControlSessionRequest}}{\mref{ControlSessionResult}}

Creates a legacy control session for a provider (v4 API).

\osubsection{connect-consumer}{\mref{ConnectConsumerRequest}}{\mref{ControlSessionResult}}

Connects a consumer to a legacy control session.

\osubsection{close-sessions}{\mref{ClosePortsRequest}}{\pref{List}$<$\allowbreak{}\pref{Map}$<$\allowbreak{}\pref{String}, \pref{Object}$>$$>$}

Closes legacy control sessions by port; returns the ports that failed.

\osubsection{open-provider-session}{\mref{ProviderSessionRequest}}{\mref{SessionEndpoint}}

Gatekeeper only: opens a relay session towards a local provider for a negotiated foreign consumer.

\osubsection{open-consumer-session}{\mref{ConsumerSessionRequest}}{\mref{SessionEndpoint}}

Gatekeeper only: opens a local listener that forwards a local consumer's requests to the provider cloud.

\newpage

\section{Information Model}
\label{sec:model}

Here, all data objects that can be part of the \texttt{gate\allowbreak{}Tunnel\allowbreak{}Connection} service are listed in alphabetic order.
Note that each subsection, which describes one type of object, begins with the \textit{struct} keyword, which is used to denote a collection of named fields, each with its own data type.
As a complement to the explicitly defined types in this section, there is also a list of implicit primitive types in Section \ref{sec:model:primitives}, which are used to represent things like hashes and identifiers.

\msubsection{struct}{ClosePortsRequest}




\noindent\begin{tabularx}{\textwidth}{| >{\raggedright\arraybackslash}p{3.9cm} | >{\raggedright\arraybackslash\small}p{5.3cm} | X |} \hline
\rowcolor{gray!33} Field & Type & Description \\ \hline
ports & \pref{List}$<$\allowbreak{}\pref{Number}$>$ & Control ports to close (a bare JSON array is also accepted). \\ \hline
\end{tabularx}



\msubsection{struct}{ConnectConsumerRequest}




\noindent\begin{tabularx}{\textwidth}{| >{\raggedright\arraybackslash}p{3.9cm} | >{\raggedright\arraybackslash\small}p{5.3cm} | X |} \hline
\rowcolor{gray!33} Field & Type & Description \\ \hline
queue\allowbreak{}Id & \pref{Identifier} & Relay queue of the provider session. \\ \hline
\end{tabularx}



\msubsection{struct}{ConsumerSessionRequest}




\noindent\begin{tabularx}{\textwidth}{| >{\raggedright\arraybackslash}p{3.9cm} | >{\raggedright\arraybackslash\small}p{5.3cm} | X |} \hline
\rowcolor{gray!33} Field & Type & Description \\ \hline
session\allowbreak{}Id & \pref{Identifier} & Session identifier from the provider cloud. \\ \hline
consumer & \pref{Name} & Consumer system. \\ \hline
provider\allowbreak{}Cloud & \pref{String} & Provider cloud. \\ \hline
provider & \pref{Name} & Provider system. \\ \hline
service\allowbreak{}Definition & \pref{Name} & Service definition. \\ \hline
peer & \pref{String} & Peer GateTunnel common name. \\ \hline
\end{tabularx}



\msubsection{struct}{ControlSessionRequest}




\noindent\begin{tabularx}{\textwidth}{| >{\raggedright\arraybackslash}p{3.9cm} | >{\raggedright\arraybackslash\small}p{5.3cm} | X |} \hline
\rowcolor{gray!33} Field & Type & Description \\ \hline
consumer & \pref{Map}$<$\allowbreak{}\pref{String}, \pref{Object}$>$ & Consumer system descriptor. \\ \hline
provider & \pref{Map}$<$\allowbreak{}\pref{String}, \pref{Object}$>$ & Provider system descriptor. \\ \hline
consumer\allowbreak{}Cloud & \pref{Map}$<$\allowbreak{}\pref{String}, \pref{Object}$>$ & Consumer cloud. \\ \hline
provider\allowbreak{}Cloud & \pref{Map}$<$\allowbreak{}\pref{String}, \pref{Object}$>$ & Provider cloud. \\ \hline
relay & \pref{Map}$<$\allowbreak{}\pref{String}, \pref{Object}$>$ & Relay descriptor. \\ \hline
service\allowbreak{}Definition & \pref{Name} & Service definition. \\ \hline
queue\allowbreak{}Id & \pref{Identifier} & Relay queue; generated if absent. \\ \hline
peer\allowbreak{}Name & \pref{String} & Peer name; generated if absent. \\ \hline
\end{tabularx}



\msubsection{struct}{ControlSessionResult}




\noindent\begin{tabularx}{\textwidth}{| >{\raggedright\arraybackslash}p{3.9cm} | >{\raggedright\arraybackslash\small}p{5.3cm} | X |} \hline
\rowcolor{gray!33} Field & Type & Description \\ \hline
session\allowbreak{}Id & \pref{Identifier} & Session identifier. \\ \hline
queue\allowbreak{}Id & \pref{Identifier} & Relay queue. \\ \hline
peer\allowbreak{}Name & \pref{String} & Peer name. \\ \hline
server\allowbreak{}Port & \pref{Number} & Local control port. \\ \hline
state & \pref{String} & Session state. \\ \hline
\end{tabularx}



\msubsection{struct}{ProviderSessionRequest}




\noindent\begin{tabularx}{\textwidth}{| >{\raggedright\arraybackslash}p{3.9cm} | >{\raggedright\arraybackslash\small}p{5.3cm} | X |} \hline
\rowcolor{gray!33} Field & Type & Description \\ \hline
session\allowbreak{}Id & \pref{Identifier} & Session identifier. \\ \hline
consumer & \pref{Name} & Consumer system. \\ \hline
consumer\allowbreak{}Cloud & \pref{Name} & Consumer cloud, must match the peer certificate. \\ \hline
provider & \pref{Name} & Provider system. \\ \hline
service\allowbreak{}Definition & \pref{Name} & Service definition. \\ \hline
peer & \pref{String} & Peer GateTunnel common name. \\ \hline
target & \pref{Map}$<$\allowbreak{}\pref{String}, \pref{Object}$>$ & Provider endpoint: \texttt{address}, \texttt{port}, \texttt{basePath} and the negotiated \texttt{operations}. \\ \hline
\end{tabularx}



\msubsection{struct}{SessionEndpoint}




\noindent\begin{tabularx}{\textwidth}{| >{\raggedright\arraybackslash}p{3.9cm} | >{\raggedright\arraybackslash\small}p{5.3cm} | X |} \hline
\rowcolor{gray!33} Field & Type & Description \\ \hline
address & \pref{String} & Advertised address of the local session listener. \\ \hline
port & \pref{Number} & Port of the session listener. \\ \hline
session\allowbreak{}Id & \pref{Identifier} & Session identifier. \\ \hline
\end{tabularx}



\subsection{Primitives}
\label{sec:model:primitives}

Types and structures mentioned throughout this document that are assumed to be available to implementations of this service.
The concrete interpretations of each of these types and structures must be provided by any IDD document claiming to implement this service.

\noindent\begin{tabularx}{\textwidth}{| p{4.6cm} | X |} \hline
\rowcolor{gray!33} Type & Description \\ \hline
\pdef{Identifier} & A string that identifies a record, e.g. a UUID or a composite identifier. \\ \hline
\pdef{List} & An \textit{array} of a known number of items, each having type A. \\ \hline
\pdef{Map} & A set of key-value pairs. \\ \hline
\pdef{Name} & A string identifier that is intended to be both human and machine-readable. \\ \hline
\pdef{Number} & An integer number. \\ \hline
\pdef{Object} & Any structured value. \\ \hline
\pdef{String} & A chain of characters. \\ \hline
\end{tabularx}

\newpage

\bibliographystyle{IEEEtran}
\bibliography{bibliography}

\newpage

\section{Revision History}
\subsection{Amendments}

\noindent\begin{tabularx}{\textwidth}{| p{1cm} | p{3cm} | p{2cm} | X | p{4cm} |} \hline
\rowcolor{gray!33} No. & Date & Version & Subject of Amendments & Author \\ \hline

1 & 2026-10-22 & \arrowversion & Document created from the Arrowhead LaTeX templates for the 5.2.2 Java Spring implementation, including the local extensions of this workspace. & Jerker Delsing \\ \hline
2 & 2026-10-30 & \arrowversion & Checked against the primary Java source tree (Eclipse-Arrowhead-java): services, operations, paths and data models unchanged. & Claude Opus 5.5 (draft) \\ \hline
\end{tabularx}

\subsection{Quality Assurance}

\noindent\begin{tabularx}{\textwidth}{| p{1cm} | p{3cm} | p{2cm} | X |} \hline
\rowcolor{gray!33} No. & Date & Version & Approved by \\ \hline

1 & & \arrowversion  &  \\ \hline

\end{tabularx}

\end{document}
